\documentclass[12pt]{article}
\usepackage[utf8]{inputenc}
\usepackage[T1]{fontenc}
\usepackage{hyperref}
\title{Wappo Language Resources for K-12 Curriculum Development}
\author{Compiled by Claude Sonnet 5 (claude-sonnet-5) for Andr\'as Kornai}
\date{}
\begin{document}
\maketitle
\section{Introduction}
Wappo (ISO 639-3 wao) is a dormant Yukian language of northern California. The last first-language speaker, Laura Fish Somersal, died in 1990. A considerable amount was documented in the 20th century.
\section{Primary Educational Resources}
No ready-made school materials were found in search results. A 2013 paper reports a plan, after a 2012 Breath of Life workshop, to produce a teaching grammar, a bilingual dictionary and youth-focused spoken instruction. I could not confirm whether these were completed.
\section{Linguistic and Reference Materials}
\begin{itemize}
\item Wappo language, overview (mentions Radin's 1924 texts and 1929 grammar, Sawyer's 1965 English-Wappo dictionary, and Thompson, Park and Li's 2006 reference grammar): \url{https://en.wikipedia.org/wiki/Wappo_language}
\item UC Berkeley Wappo page: \url{https://cla.berkeley.edu/languages/wappo.html}
\item Endangered Languages Project bibliography: \url{https://endangeredlanguages.com/lang/1150/bibliography?hl=es}
\end{itemize}
\section{Religious and Community Resources}
None found in search results.
\section{Digital and Audio Resources}
Sawyer's 1960s recordings with Somersal are described in the overview above; I did not find a direct archive link.
\section{Curriculum Development Context}
\begin{itemize}
\item Marks and Harp, ``Language revitalization in Northern California: Awakening the Wappo language'': \url{https://scholarspace.manoa.hawaii.edu/items/13c4e2d1-ed84-4b02-8672-a4bc9557464a}
\end{itemize}
Wappo is not known to be taught in schools; a curriculum would be built on archival documentation.
\end{document}